%%Ha_Hu
T est converti en Kelvin, si besoin est, puis :
\[
\text{Ha} = 1000  \frac{\frac{\text{Hu}}{100} . e_{w}(T)}{ R_v . T}
\]
avec :\newline
$R_v$ = 461.51  J kg$^{-1}$ K$^{-1}$ (constante sp�cifique de la vapeur d'eau).\newline
La fonction $e_w(T)$ (tension de vapeur saturante au-dessus de l'eau liquide)
est celle de la biblioth�que thermodynamique d'AIDA, d�crite en annexe 
\ref{annexe-thermo-ew}.
